\printedpage{24}
% head: A. Grothendieck : Dérivateurs, Ch. V, Édition des Universités de Montpellier II et Paris VII
% foot: Version 9 septembre 2001

\[
\begin{tikzcd}
\Hom_C(X,K) \arrow[r,"\alpha_X"] \arrow[dr,dashed] & \Hom_{C'}(\varphi X,\varphi K) \arrow[d,dash,"\wr"',"\mathrm{adj}"] \\
 & \Hom_C(X,\psi\varphi K)
\end{tikzcd}
\]
is deduced from the adjunction arrow
\[
K \longrightarrow \psi\varphi K
\]
by taking $\Hom_C(X,-)$, so $\alpha_X$ is an isomorphism for every $X$ if and only if this adjunction arrow is an isomorphism. So
\[
\text{(ii)} \Longleftrightarrow \text{(ii$'$)} \Longleftrightarrow \text{(ii$''$)}.
\]
On the other hand,
\[
\text{(ii$''$)} \Longrightarrow \text{(i)} \Longrightarrow \text{(i$'$)} \Longrightarrow \text{(i$''$)}
\]
are tautological, bearing in mind for (i) $\Longrightarrow$ (i) \texttt{[probably (i) $\Longrightarrow$ (i$'$)]} that the adjunction arrows
\[
X \xrightarrow{u} \psi\varphi X \qquad\qquad\qquad \text{in } C
\]
are indeed in $W$, for we have a commutative triangle
\[
\begin{tikzcd}[column sep=large]
\varphi(X) \arrow[r,"\varphi(u)"] \arrow[dr,"\id_{\varphi(X)}"] & \varphi(\psi\varphi X) = (\varphi\psi)(\varphi X) \arrow[d,"{\begin{array}{l}\wr\ \mathrm{adj}\ \varphi\psi \longrightarrow 1,\ \text{which is an}\\ \text{isomorphism since } \psi \text{ is fully faithful}\end{array}}"] \\
 & \varphi(X),
\end{tikzcd}
\]
so $\varphi(u)$ is indeed an isomorphism.

It remains to prove (i$''$) $\Longrightarrow$ (ii$'$), i.e. we suppose

\noindent\msp{page 39}%
that the functor $X \longrightarrow \Hom(X,K)$ transforms the particular adjunction morphism
\[
K \xrightarrow{i} \psi\varphi K
\]
into
\[
\begin{aligned}
u \quad &\longrightarrow \quad u\circ i\\
\Hom(\psi\varphi(K),K) \quad &\longrightarrow \quad \Hom(K,K),
\end{aligned}
\]
which is a bijection. So $\id_K$ is in the image, let $p$ be an element of the first member which gives rise to it

\printedpage{25}
% head: A. Grothendieck : Dérivateurs, Ch. V, Édition des Universités de Montpellier II et Paris VII
% foot: Version 9 septembre 2001

\[
\begin{tikzcd}
K \arrow[r,hook,shift left,"i"] & \psi\,\underbrace{\varphi(K)}_{K'} \arrow[l,shift left,"p"]
\end{tikzcd}
\qquad\qquad {\scriptstyle \text{so } p\circ i = \id_K}
\]
Let us consider
\[
ip = q : \psi(K') \longrightarrow \psi(K') \qquad\qquad\qquad (\text{NB } q^2 = q),
\]
since $\psi$ is fully faithful, there exists a unique
\[
q_0 : K' \longrightarrow K' \qquad\qquad\qquad (K' \overset{\text{def}}{=} \varphi(K))
\]
such that we have
\[
q = \psi(q_0).
\]
The relation $pi = \id_K$ gives, on applying $\varphi$,
\[
\varphi(p)\varphi(i) = \id,
\]
so since $\varphi(i)$ is an isomorphism (it is
\[
\begin{tikzcd}[column sep=large]
\varphi(K) \arrow[r,"\sim"] \arrow[dr,"\sim"{sloped,pos=0.45},"\id"{pos=0.6}] & \varphi(\psi\varphi)(K) = (\varphi\psi)(\varphi K) \arrow[d,"\wr"',"\mathrm{adj}"] \\
 & \varphi(K),
\end{tikzcd}
\]
making the triangle opposite commutative), $\varphi(p)$ is one. So $\varphi(q) = \varphi(i)\varphi(p)$ is one, since $\varphi(i)$ and $\varphi(i)$ \texttt{[rather $\varphi(p)$]} are. But
\[
\varphi(q) = \varphi\psi(q_0) \quad \underbrace{\xrightarrow{\ \sim\ }}_{\mathrm{adj}\ \varphi\psi\ \xrightarrow{\sim}\ \id} \quad q_0,
\]
so $q_0$ is an isomorphism, hence also $q = \psi q_0$. Since $q^2 = q$, this implies $q = \id$. So $i$ and $p$ are isomorphisms inverse to one another, which establishes (ii$''$) \texttt{[rather (ii$'$)]} and completes the proof.

\textbf{Corollary:} \textit{Let us consider functors $C^o \xrightarrow{H} \Ens$, and $C'^o \xrightarrow{H'} \Ens$, such that $H \simeq H'\circ\varphi$. For $H$ to be representable, it is necessary and sufficient that $H'$ be so, and in this case, the object $K$ of $C$}

\noindent\msp{page 40}%
\textit{which represents $H$ satisfies the equivalent conditions of the proposition. If $K$, $K'$ represent them, we have canonical isomorphisms $K' \simeq \varphi K$, $K \simeq \psi K'$.}

If $H'$ is representable by $K' \in \Ob C'$, we will therefore have
\[
H'\circ\varphi(X) \simeq \Hom_{C'}(\varphi X,K') \overset{\mathrm{adj}}{\simeq} \Hom_C(X,\psi(K'))
\]
\printedpage{26}
% head: A. Grothendieck : Dérivateurs, Ch. V, Édition des Universités de Montpellier II et Paris VII
% foot: Version 9 septembre 2001
so $H$ is representable by $K = \psi(K')$, satisfying condition (ii) of the proposition. Conversely, if $H$ is representable by an object $K$, then $K$ satisfies condition (i), so by (ii$''$) we have
\[
\Hom_C(X,K) \xrightarrow{\sim} \Hom_{C'}(\varphi X,\varphi K),
\]
so the functor $H' : C^o \longrightarrow \Ens$ which factors $H$ is given explicitly as $X' \longmapsto \Hom_{C'}(X',\varphi K)$, hence is representable by $\varphi K$. % [sic] C^o -> C'^o, H' is a functor on C'^o (cf. the Corollaire)

Let us pass to ind-objects. The functors $\varphi$, $\psi$ pass to ind-objects, whence functors
\[
\begin{tikzcd}[column sep=large]
\Ind C \arrow[r,shift left,"\Psi = \Ind\psi"] \arrow[r,shift right,"\Phi = \Ind\varphi"'] & \Ind C'.
\end{tikzcd} % [sic] both arrows printed pointing to Ind C' -> the arrow Psi = Ind psi apparently intended from Ind C' to Ind C, psi goes from C' to C
\]
We will again have the adjunction formula for $X = (X_i)_{i\in I}$ in $\Ind C$, $X' = (X'_{i'})_{i'\in I'}$ in $\Ind C'$, for
\[
\begin{aligned}
\Hom(X,\Psi X') &= \varprojlim_i \varinjlim_{i'} \Hom(X_i,\psi(X'_{i'}))\\
&= \varprojlim_i \varinjlim_{i'} \Hom(\varphi X_i,X'_{i'})\\
&= \Hom(\Phi X,X').
\end{aligned}
\]

\noindent\msp{page 41}%
We prove in the same way that $\Psi$ is fully faithful, for if $X' = (X'_i)_{i\in I}$, $Y' = (Y'_j)_{j\in J}$ are objects of $\Ind C'$, we will have
\[
\begin{tikzcd}[column sep=small]
\Hom(X',Y') \arrow[d,"\Psi_{X',Y'}"] & \simeq \varprojlim\limits_i \varinjlim\limits_j & \Hom_{C'}(X'_i,Y'_j) \arrow[d,"\psi_{X'_i,Y'_j}"] \\
\Hom(\Psi X',\Psi Y') & \simeq \varprojlim\limits_i \varinjlim\limits_j & \Hom_C(\psi(X'_i),\psi(Y'_j)),
\end{tikzcd}
\]
so $\Psi_{X',Y'}$ is an isomorphism if the $\psi_{X'_i,Y'_j}$ are.

So proposition 1 can be applied to the situation $\Phi$, $\Psi$ between ind-objects. I am going to spell it out in the most convenient form:

\textbf{Proposition 2.} \textit{Let $K = (K_i)_{i\in I}$ be an ind-object of $C$, and let $H$ be the functor}
\[
\left\{
\begin{aligned}
H : X \quad &\longmapsto \quad \underbrace{\Hom_C(X,K)}_{K(X)} \overset{\text{def}}{=} \varinjlim_i \Hom_C(X,K_i)\\
C^o \quad &\longrightarrow \quad \Ens
\end{aligned}
\right.
\]
\textit{which it ind-represents. Equivalent conditions}

\begin{enumerate}[label={}]
\item[(i)] \textit{The functor $H$ factors (up to isomorphism) through $C'$ via $\varphi$, i.e. transforms the $u \in W_\varphi$ into isomorphisms.}
\printedpage{27}
% head: A. Grothendieck : Dérivateurs, Ch. V, Édition des Universités de Montpellier II et Paris VII
% foot: Version 9 septembre 2001
\item[(i$'$)] \textit{Same condition for the adjunction arrows $X \longrightarrow \psi\varphi(X)$ (which are in $W$).}
\item[(i$''$)] \textit{Same condition for the adjunction arrow $K \longrightarrow \psi\varphi(K)$.}
\item[(ii)] \textit{$K$ belongs to the essential image of} $\Ind C'$ \textit{under $\Psi$.}
\item[(ii$'$)] $K \longrightarrow \Psi\Phi(K)$ \textit{is an} isomorphism \textit{in} $\Ind C$.) % [sic] ".)" -> ".", closing parenthesis without opening parenthesis
\item[(ii$''$)] \textit{The functorial homomorphism}
\[
K(X) = \Hom_C(X,K) \longrightarrow \Hom_{C'}(\varphi(X),\varphi(K))
\]
\texttt{[rather} $\Hom_{\mathrm{Ind}(C)}(X,K) \longrightarrow \Hom_{\mathrm{Ind}(C')}(\Phi(X),\Phi(K))$\texttt{, identifying the object} $X$ \texttt{of} $C$ \texttt{with its image in} $\mathrm{Ind}(C)$\texttt{]} \textit{is an isomorphism.}
\end{enumerate}

\textsc{Proof.} The equivalence of (ii), (ii$'$), (ii$''$) is proposition 1 applied to $\Phi$ and $\Psi$, except that we must note that (ii$''$) obviously remains valid if we replace $X$ in it by an arbitrary ind-object $(X_j)_{j\in J}$. We have

\noindent\msp{page 42}%
\[
\text{(ii$''$)} \Longrightarrow \text{(i)} \Longrightarrow \text{(i$'$)}
\]
tautologically, it remains to prove that (i$'$) implies (i$''$) and (i$''$) (ii$'$). We cannot apply proposition 1 directly, for hypothesis (i$'$) here concerns only the objects of $\Ind C$ which come from $C$, and it is not obvious a priori that it passes to arbitrary objects. But it is in the case of a homomorphism of ind-objects $(X_i)$, $(Y_i)$ on the same index set, given by an ordinary (term-by-term) homomorphism of direct systems on $I$, $X_i \longrightarrow Y_i$: if the $X_i \longrightarrow Y_i$ are in $W$, hypothesis (i) implies that $H(Y) \longrightarrow H(X)$ is an isomorphism, and likewise for (i$'$) for $\underbrace{X}_{=(X_i)_i} \;\longrightarrow\; \underbrace{\Psi\Phi(X)}_{=(\psi\varphi(X_i))_i}$\,, where the adjunction homomorphism is given term by term by the adjunction homomorphisms $X_i \longrightarrow \psi\varphi(X_i)$. So hypothesis (i$'$) here implies hypothesis (i$'$) of proposition 1 for $\Phi$ and $\Psi$, hence also (ii$'$) (by proposition 1), and we win. % [sic] "et (i'') (ii')" -> "et (i'') => (ii')", the implication sign or verb is missing

Let us return to the situation where $C = \Cathot$, $C' = \Hot$, $\varphi$ the functor \texttt{[of]} localization. If $K \in \Ob\Cathot$, then $\psi\varphi(K)$ is called \textit{the Thomason-Kan envelope} of $K$ - if $K$ is a category, its envelope $\bar{K}$ and the arrow $K \longrightarrow \bar{K}$ are only defined in the category Cat\,hot, i.e. $\bar{K}$ is only defined ``up to homotopism''. % reading resolved: no underline is printed under psi phi(K).

\noindent\msp{page 43}%
But we can always \textit{choose} an object $\bar{K}$ of Cat, and an arrow
\[
i : K \longrightarrow \bar{K} \qquad\qquad\qquad \textit{in } \Cat,
\]
\printedpage{28}
% head: A. Grothendieck : Dérivateurs, Ch. V, Édition des Universités de Montpellier II et Paris VII
% foot: Version 9 septembre 2001
not only in Cat\,hot, in such a way moreover that $i$ is a ``Thomason \textit{cofibration}'' (and in particular an open immersion?). It follows that if we have a direct system
\[
K = (K_n)_{n\in\mathbf{N}}
\]
indexed by $\mathbf{N}$, and if we choose for each $n$ a Thomason-Kan envelope
\[
K_n \overset{\alpha_n}{\hookrightarrow} \bar{K}_n
\]
in this way, we can find transition functors
\[
\bar{K}_n \xrightarrow{\bar{i}_n} \bar{K}_{n+1}
\]
such that the squares
\[
\begin{tikzcd}
K_n \arrow[r,hook,"\alpha_n"] \arrow[d,"i_n"'] & \bar{K}_n \arrow[d,"\bar{i}_n"] \\
K_{n+1} \arrow[r,hook,"\alpha_{n+1}"] & \bar{K}_{n+1}
\end{tikzcd}
\]
are commutative, so that we thus have a homomorphism of \textit{direct systems of} Cat
\[
K \longrightarrow \bar{K}
\]
and a fortiori a homomorphism of ind-objects of Cat. This being so, conditions (i), (ii$'$) of proposition 2 show that the functor ind-represented by $K$ \textit{in} $C = \Cathot$ transforms ordinary

\noindent\msp{page 44}%
quasi-isomorphisms (weak equivalences) into bijections, \textit{if and only if} the preceding homomorphism is an isomorphism of ind-objects of Ind\,Cat\,hot.

So that will be the natural condition to impose on the ind-categories of paths
\[
\Chh_{\Sigma}(\mathcal{M};a,b) \qquad\qquad\qquad (a,b \in \mathrm{ob}\,\mathcal{M})
\]
coming from a ``category of models'' $\mathcal{M}$, $\Sigma$.

Of course, having at our disposal an explicit direct system of categories (not only of objects of Cat\,hot)
\[
K_n = \underline{\Ch}_{\Sigma,n}(\mathcal{M};a,b),
\]
it is tempting to pass to the direct limit
\[
K_\infty = \varinjlim_n \underline{\Ch}_{\Sigma,n}(\mathcal{M};a,b) \overset{\text{def}}{=} \underline{\Ch}_{\Sigma}(\mathcal{M};a,b),
\]
and this all the more so since the transition functors are monomorphisms and fully faithful, so that the $K_n$ are identified with a sequence of full subcategories of